\documentclass[12pt,a4paper]{article}
% Drake_preamble.tex - LaTeX Preamble Template
% 作者: 謝慕揚, MD, PhD, FESC
% 
% 使用說明:
% 1. 表格請使用 longtable，不要有垂直分隔線 |
% 2. 表格內換行使用 \newline，不要使用 \\
% 3. 藥物名稱、檢驗、檢查請維持英文
% 4. Reference 格式請使用 \href 加上驗證過的 DOI

% Including essential packages for Chinese typesetting
\usepackage{xeCJK}
\usepackage{fontspec}
% Configuring Chinese font with Noto Serif CJK TC, fallback to SimSun
\IfFontExistsTF{Noto Serif CJK TC}{
  \setCJKmainfont{Noto Serif CJK TC}
  \setCJKsansfont{Noto Serif CJK TC}
  \setCJKmonofont{Noto Serif CJK TC}
}{\setCJKmainfont{SimSun}
  \setCJKsansfont{SimSun}
  \setCJKmonofont{SimSun}
}

% Including other necessary packages
\usepackage{geometry}
\usepackage{titlesec}
\usepackage{enumitem}
\usepackage{xcolor}
\usepackage{colortbl}
\usepackage{tcolorbox}
\usepackage{booktabs}
\usepackage{longtable}
\usepackage{array}
\usepackage{graphicx}
\usepackage{tikz}
\usepackage{fancyhdr}
\usepackage{multicol}
\usepackage{datetime2}
\usepackage{hyperref}
\usepackage{mdframed}
\usepackage{amsmath}
\usepackage{amssymb}
\usepackage{multirow}

\hypersetup{
    colorlinks=true,
    linkcolor=emergency_blue,
    citecolor=emergency_blue,
    urlcolor=emergency_blue,
    pdfborder={0 0 0}
}

% Defining page geometry
\geometry{
  top=2.5cm,
  bottom=2.5cm,
  left=2cm,
  right=2cm,
  headheight=25pt
}

% Adjusting topmargin to compensate for increased headheight
\addtolength{\topmargin}{-10pt}

% Defining colors
\definecolor{emergency_red}{RGB}{186,24,27}
\definecolor{emergency_blue}{RGB}{0,114,188}
\definecolor{emergency_gray}{RGB}{120,120,120}
\definecolor{highlight_yellow}{RGB}{255,250,205}

% Setting title formats
\titleformat{\section}[block]{\Large\bfseries\color{emergency_red}}{\thesection}{10pt}{}
\titleformat{\subsection}[block]{\large\bfseries\color{emergency_blue}}{\thesubsection}{8pt}{}
\titleformat{\subsubsection}[block]{\normalsize\bfseries\color{emergency_gray}}{\thesubsubsection}{6pt}{}

% Defining custom environment for important notes
\newmdenv[
    linecolor=emergency_red,
    backgroundcolor=highlight_yellow,
    linewidth=2pt,
    topline=true,
    bottomline=true,
    leftline=true,
    rightline=true,
    innertopmargin=10pt,
    innerbottommargin=10pt,
    innerleftmargin=10pt,
    innerrightmargin=10pt
]{importantbox}

% Defining note command with tcolorbox
\newcommand{\note}[1]{
    \par
    \noindent
    \begin{tcolorbox}[
        colback=emergency_blue!5!white,
        colframe=emergency_blue,
        title=\textbf{重點提醒},
        fonttitle=\bfseries,
        boxrule=0.5pt,
        arc=3pt,
        left=6pt,
        right=6pt,
        top=6pt,
        bottom=6pt
    ]
    #1
    \end{tcolorbox}
}

% Key findings box
\newcommand{\keyfinding}[1]{
    \par
    \noindent
    \begin{tcolorbox}[
        colback=yellow!10!white,
        colframe=orange!75!black,
        title=\textbf{核心發現摘要},
        fonttitle=\bfseries,
        boxrule=0.5pt,
        arc=3pt
    ]
    #1
    \end{tcolorbox}
}

% Defining custom date format for ISO 8601
\DTMnewdatestyle{iso}{
  \renewcommand{\DTMdisplaydate}[4]{\number##1-\number##2-\number##3}
  \renewcommand{\DTMDisplaydate}[4]{\number##1-\number##2-\number##3}
}
\DTMsetdatestyle{iso}
\newcommand{\mydate}{\DTMtoday}


\pagestyle{fancy}
\fancyhf{}
\fancyhead[L]{\textcolor{emergency_red}{European Heart Journal 教學講義}}
\fancyhead[R]{\textcolor{emergency_blue}{心血管預防專題}}
\fancyfoot[C]{\textcolor{emergency_gray}{\thepage}}
\renewcommand{\headrulewidth}{1pt}
\renewcommand{\headrule}{\hbox to\headwidth{\color{emergency_red}\leaders\hrule height \headrulewidth\hfill}}

\begin{document}

\begin{center}
{\LARGE\bfseries\textcolor{emergency_red}{European Heart Journal}}\\[0.5cm]
{\Large\textbf{Improving Prevention: Traditional and Non-Traditional Risk Factors}}\\[0.3cm]
{\large 改善心血管預防：傳統與非傳統危險因子}\\[0.5cm]
{\normalsize \href{https://doi.org/10.1093/eurheartj/ehaf1092}{Eur Heart J 2026;47:279-284. DOI: 10.1093/eurheartj/ehaf1092}}\\[0.3cm]
{\normalsize 整理：謝慕揚 MD, PhD, FESC \quad 日期：\mydate}
\end{center}

\vspace{0.5cm}

\begin{tcolorbox}[colback=yellow!10!white,colframe=orange!75!black,title=\textbf{本期核心重點}]
\begin{enumerate}
    \item \textbf{SMuRF-less but inflamed：}無傳統危險因子但hsCRP升高的女性，30年CV事件風險仍顯著增加，Statin可降低38\%風險
    \item \textbf{懷孕期間Statin安全性：}挪威全國性研究顯示，第一孕期statin暴露與先天畸形無顯著相關
    \item \textbf{年輕人降壓治療：}義大利真實世界研究證實，18-39歲高血壓患者adherence可降低22\% CV住院風險
    \item \textbf{Troponin與失智：}中年subclinical myocardial injury與25年後失智風險增加10\%相關
\end{enumerate}
\end{tcolorbox}

\tableofcontents
\newpage

%============================================================
\section{State of the Art Review：社經弱勢與指引落實的障礙}
%============================================================

\subsection{重要發現}

\begin{itemize}
    \item \textbf{社經弱勢族群}（low socioeconomic position）有較高CVD負擔與死亡率
    \item 令人擔憂的是：這些族群同時獲得\textbf{較差的guideline-recommended care}
    \item 形成雙重弱勢：疾病風險高、醫療照護品質低
\end{itemize}

\subsection{建議對策}

\begin{longtable}{p{4cm}p{10cm}}
\toprule
\textbf{層面} & \textbf{建議} \\
\midrule
\endfirsthead
\bottomrule
\endfoot

政府層級 & 降低心血管健康不平等、改善人口教育 \\
\midrule
醫療體系 & 優先處理可修正的照護障礙 \\
\midrule
研究層級 & RCT設計應納入社經弱勢族群 \\
\midrule
臨床實務 & 使用quality indicators監測照護品質\newline
調整risk scores納入社經因素 \\
\bottomrule
\end{longtable}

%============================================================
\section{Fast Track：hsCRP與SMuRF-less女性的心血管風險}
%============================================================

\subsection{研究背景}

\begin{itemize}
    \item 傳統預防聚焦於四大\textbf{SMuRFs}（Standard Modifiable Risk Factors）：
    \begin{itemize}
        \item 高血壓
        \item 血脂異常
        \item 糖尿病
        \item 吸菸
    \end{itemize}
    \item 但大量CV事件發生於\textbf{無SMuRF}的個體，尤其是\textbf{女性}
    \item \textbf{hsCRP}是敏感的發炎指標，但在SMuRF-less族群的角色不明
\end{itemize}

\subsection{Women's Health Study設計}

\begin{longtable}{p{4cm}p{10cm}}
\toprule
\textbf{項目} & \textbf{內容} \\
\midrule
\endfirsthead
\bottomrule
\endfoot

研究族群 & 12,530位初始健康、無SMuRF的女性 \\
\midrule
追蹤時間 & 30年 \\
\midrule
基線檢測 & hsCRP \\
\midrule
主要終點 & Incident CHD、ischemic stroke、total CV events \\
\midrule
分析方法 & Quintile分析、spline regression \\
\bottomrule
\end{longtable}

\subsection{主要結果}

\keyfinding{
30年追蹤中，共973例首次重大CV事件。發生事件者的基線hsCRP顯著較高（中位數2.22 vs 1.50 mg/L，P$<$0.0001）。
}

\begin{longtable}{p{5cm}ccc}
\toprule
\textbf{終點} & \textbf{HR (Q5 vs Q1)} & \textbf{P-trend} & \textbf{每quintile增加HR} \\
\midrule
\endfirsthead
\bottomrule
\endfoot

CHD & 2.23 & $<$0.0001 & 1.21 \\
Ischemic stroke & 1.69 & $<$0.0001 & -- \\
Total CV events & 1.74 & $<$0.0001 & -- \\
\bottomrule
\end{longtable}

\textbf{以臨床常用閾值分析（hsCRP $>$3 vs $<$1 mg/L）：}
\begin{itemize}
    \item CHD風險增加\textbf{77\%}
    \item Ischemic stroke風險增加\textbf{39\%}
    \item Total CV events風險增加\textbf{52\%}
\end{itemize}

\note{
\textbf{Editorial觀點（Arnold \& Koenig）：}\\
「SMuRF-less but inflamed」表型挑戰現有預防概念。我們需要重繪預防調色盤：繼續尋找SMuRFs的熟悉藍色，但也要學會辨識低度發炎的隱藏紅色——在它發作前就發出危險信號。\\[0.3cm]
\textcolor{red}{\textbf{重要：}}Statin治療可在此族群降低\textbf{38\%}風險！
}

%============================================================
\section{Clinical Research：懷孕期間Statin使用與先天畸形}
%============================================================

\subsection{研究背景}

\begin{itemize}
    \item Statins傳統上在懷孕期間\textbf{禁忌使用}
    \item 主要顧慮：可能對胎兒有害
    \item 但實際風險數據有限
\end{itemize}

\subsection{挪威全國性研究設計}

\begin{longtable}{p{4cm}p{10cm}}
\toprule
\textbf{項目} & \textbf{內容} \\
\midrule
\endfirsthead
\bottomrule
\endfoot

資料來源 & 挪威全國登錄資料（2005-2018年） \\
\midrule
研究對象 & 所有懷孕女性 \\
\midrule
暴露定義 & 第一孕期statin處方 \\
\midrule
對照組 & Statin discontinuers（懷孕前停藥） \\
\midrule
主要終點 & Any congenital malformation\newline
Major malformation\newline
Minor malformation \\
\bottomrule
\end{longtable}

\subsection{主要結果}

\begin{longtable}{lccc}
\toprule
\textbf{族群} & \textbf{人數} & \textbf{畸形率} \\
\midrule
\endfirsthead
\bottomrule
\endfoot

Statin未暴露 & $>$800,000 & 4.3\% \\
Statin停藥者 & $>$1,255 & 5.9\% \\
Statin暴露 & 283 & 6.7\% \\
\bottomrule
\end{longtable}

\textbf{調整後Odds Ratios（暴露 vs 未暴露）：}
\begin{itemize}
    \item Any malformation: OR 1.30（非顯著）
    \item Major malformation: OR 1.15（非顯著）
    \item Minor malformation: OR 1.47（非顯著）
\end{itemize}

\textbf{暴露 vs Discontinuers：}
\begin{itemize}
    \item Any malformation: OR 1.01
    \item Major malformation: OR 1.08
    \item Minor malformation: OR 0.94
\end{itemize}

\begin{tcolorbox}[colback=emergency_blue!5!white,colframe=emergency_blue,title=\textbf{臨床意義}]
\begin{itemize}
    \item 本研究\textbf{未發現}第一孕期statin暴露與先天畸形的顯著關聯
    \item 不支持statin與先天畸形之間存在強烈或獨立的因果關係
    \item 對於高風險女性（如familial hypercholesterolemia），懷孕期間LLT的benefit可能outweigh risk
    \item 有助於informed shared decision-making
\end{itemize}
\end{tcolorbox}

%============================================================
\section{Clinical Research：血壓、血漿蛋白與心血管疾病}
%============================================================

\subsection{研究目的}

透過\textbf{Mendelian Randomization (MR)}系統性評估血漿蛋白與血壓的因果關係，並探討其與CVD的連結。

\subsection{主要發現}

\begin{itemize}
    \item 分析2,007種血漿蛋白
    \item 242種蛋白與BP相關
    \item 其中48種同時與CAD或stroke相關
    \item \textbf{4種蛋白}通過co-localization分析：
    \begin{itemize}
        \item \textbf{ACOX1}
        \item \textbf{FGF5}
        \item \textbf{FURIN}
        \item \textbf{MST1}
    \end{itemize}
\end{itemize}

\subsection{血壓中介效應}

Network MR顯示，這些蛋白對CAD和stroke的效應，有\textbf{30.5-77.2\%}是透過BP regulation中介。

\note{
\textbf{意義：}這項研究辨識出具有因果角色的關鍵血漿蛋白，提供高血壓相關CVD分子機轉的新見解，並指出未來治療標靶的方向。
}

%============================================================
\section{Clinical Research：年輕人降壓治療與心血管風險}
%============================================================

\subsection{研究背景}

\begin{itemize}
    \item 降壓治療在老年及中年人的效益已有充分證據
    \item 但在\textbf{年輕人（$<$40歲）}的效益從未被直接驗證
\end{itemize}

\subsection{義大利Lombardy世代研究}

\begin{longtable}{p{4cm}p{10cm}}
\toprule
\textbf{項目} & \textbf{內容} \\
\midrule
\endfirsthead
\bottomrule
\endfoot

資料來源 & Lombardy地區健保資料庫 \\
\midrule
研究對象 & $>$286,000位首次處方降壓藥物者\newline
年齡18-55歲（2009-2017年） \\
\midrule
Adherence定義 & 處方覆蓋$\geq$80\% vs $<$80\%追蹤期間 \\
\midrule
主要終點 & CV住院 \\
\midrule
次要終點 & CV死亡、全因死亡 \\
\bottomrule
\end{longtable}

\subsection{主要結果}

\begin{longtable}{lcc}
\toprule
\textbf{終點} & \textbf{HR (18-39歲)} & \textbf{HR (40-55歲)} \\
\midrule
\endfirsthead
\bottomrule
\endfoot

CV住院 & 0.78 & 0.80 \\
CV死亡 & NS & 顯著降低 \\
全因死亡 & NS & 顯著降低 \\
\bottomrule
\end{longtable}

\keyfinding{
在真實世界中，adherence to降壓治療可在年輕人（18-39歲）達到與中年人（40-55歲）\textbf{同等程度}的CV住院風險降低。即使目前兩組的adherence都偏低，年輕患者的治療教育仍應被強化。
}

%============================================================
\section{Clinical Research：心肌Troponin I與失智風險}
%============================================================

\subsection{Whitehall II Study設計}

\begin{longtable}{p{4cm}p{10cm}}
\toprule
\textbf{項目} & \textbf{內容} \\
\midrule
\endfirsthead
\bottomrule
\endfoot

研究族群 & $\sim$6,000位參與者，基線年齡45-69歲 \\
\midrule
基線檢測 & High-sensitivity cardiac troponin I（1997-99年） \\
\midrule
追蹤至 & 2023年3月（中位數25年） \\
\midrule
主要終點 & Incident dementia \\
\midrule
其他評估 & 認知測試（6波）、MRI神經影像（2012-16年） \\
\bottomrule
\end{longtable}

\subsection{主要發現}

\begin{itemize}
    \item 606例（10.1\%）失智症
    \item \textbf{Cardiac troponin每倍增}，失智風險增加\textbf{10\%}（95\% CI 3-17\%）
    \item 失智者在診斷前7-25年就有較高troponin
    \item 基線troponin $>$5.2 ng/L者，15年後：
    \begin{itemize}
        \item Grey matter volume較低（相當於老化2.7年）
        \item Hippocampal atrophy較嚴重（相當於老化3年）
    \end{itemize}
\end{itemize}

\note{
\textbf{心腦連結：}中年的subclinical myocardial injury與晚年失智風險相關。High-sensitivity troponin不僅可評估心臟健康，也可能成為腦部健康的指標。
}

%============================================================
\section{Rapid Communication：女性健美選手死亡率}
%============================================================

\subsection{研究發現}

\begin{longtable}{p{5cm}p{9cm}}
\toprule
\textbf{項目} & \textbf{數據} \\
\midrule
\endfirsthead
\bottomrule
\endfoot

研究族群 & $\sim$10,000位IFBB女性選手（2005-2020年） \\
\midrule
死亡人數 & 32例 \\
\midrule
SCD發生率 & 10.47 per 100,000 athlete-years \\
\midrule
職業選手SCD & \textcolor{red}{53.98 per 100,000 AY}（顯著較高） \\
\midrule
自殺/他殺 & 佔所有死亡$\sim$13\%（明顯高於男性） \\
\bottomrule
\end{longtable}

\begin{tcolorbox}[colback=emergency_blue!5!white,colframe=emergency_blue,title=\textbf{臨床建議}]
女性健美選手，尤其職業等級，面臨顯著的心血管與心理風險。建議：
\begin{enumerate}
    \item 加強心血管篩檢
    \item 規範performance-enhancing drugs使用
    \item 提供量身訂做的心理健康支持
\end{enumerate}
\end{tcolorbox}

%============================================================
\section{縮寫對照表}
%============================================================

\begin{longtable}{p{4cm}p{10cm}}
\toprule
\textbf{縮寫} & \textbf{全名} \\
\midrule
\endfirsthead
\bottomrule
\endfoot

SMuRF & Standard Modifiable Risk Factors (標準可修正危險因子) \\
hsCRP & High-Sensitivity C-Reactive Protein (高敏感度C反應蛋白) \\
CHD & Coronary Heart Disease (冠狀動脈心臟病) \\
LLT & Lipid-Lowering Therapy (降血脂治療) \\
MR & Mendelian Randomization (孟德爾隨機化) \\
CAD & Coronary Artery Disease (冠狀動脈疾病) \\
BP & Blood Pressure (血壓) \\
CV & Cardiovascular (心血管) \\
SCD & Sudden Cardiac Death (心因性猝死) \\
IFBB & International Fitness and Bodybuilding Federation \\
AY & Athlete-Years (運動員年) \\
FH & Familial Hypercholesterolemia (家族性高膽固醇血症) \\
\bottomrule
\end{longtable}

%============================================================
\section{參考文獻}
%============================================================

\begin{enumerate}
    \item Crea F. Improving prevention: traditional and non-traditional risk factors. \href{https://doi.org/10.1093/eurheartj/ehaf1092}{\textit{Eur Heart J}. 2026;47:279-284.}

    \item Ridker PM, et al. C-reactive protein and cardiovascular risk among women with no standard modifiable risk factors. \href{https://doi.org/10.1093/eurheartj/ehaf658}{\textit{Eur Heart J}. 2026;47:306-314.}

    \item Christensen JJ, et al. Statin use in pregnancy and risk of congenital malformations: a Norwegian nationwide study. \href{https://doi.org/10.1093/eurheartj/ehaf592}{\textit{Eur Heart J}. 2026;47:318-327.}

    \item Meena D, et al. Blood pressure, plasma proteins, and cardiovascular diseases: a network Mendelian randomization and observational study. \href{https://doi.org/10.1093/eurheartj/ehaf725}{\textit{Eur Heart J}. 2026;47:331-342.}

    \item Rea F, et al. Antihypertensive treatment in young adults and cardiovascular risk. \href{https://doi.org/10.1093/eurheartj/ehaf744}{\textit{Eur Heart J}. 2026;47:346-355.}

    \item Chen Y, et al. High-sensitivity cardiac troponin I and risk of dementia: the 25-year longitudinal Whitehall II study. \href{https://doi.org/10.1093/eurheartj/ehaf834}{\textit{Eur Heart J}. 2026;47:359-369.}

    \item Vecchiato M, et al. Mortality in female bodybuilding athletes. \href{https://doi.org/10.1093/eurheartj/ehaf789}{\textit{Eur Heart J}. 2026;47:373-375.}
\end{enumerate}

\end{document}
